\emph{(i).} By Theorem \claimref{dark-energy.frozen-point-makes-stelle-west-vector}(i) the frozen point is a vector $I$ of the tractor space, of signature $(2,4)$, with $Q(I) = \Lambda/3$. If $Q(I) \neq 0$ the space is the orthogonal sum of the line of $I$ and $I^{\perp}$, so by Fact \claimref{std-algebra.sylvester-inertia} the signature of $I^{\perp}$ is $(2,4)$ minus that of the line: $(1,4)$ for $Q(I) > 0$ and $(2,3)$ for $Q(I) < 0$. If $Q(I) = 0$, $I \in I^{\perp}$ spans its radical, and $I^{\perp}/\langle I\rangle$ is nondegenerate of signature $(1,3)$. The check builds the Gram matrix of an exact basis of $I^{\perp}$ from geometric products in $\Cl(2,4)$ and diagonalises it by congruence, for basis and non-basis $I$ of each sign of $Q(I)$, and confirms the three signatures and the residue $1 - 4 \equiv 5$.

\emph{(ii).} Let $W$ be a nondegenerate subspace of signature $(a,b)$ of a quadratic space $U$ with nondegenerate part $(P,Q)$ and radical $R$. Choose an orthogonal basis of $U$ as in Fact \claimref{std-algebra.sylvester-inertia} and let $N$ be the span of its $Q$ negative and its radical vectors, on which the form is negative semidefinite. A positive-definite subspace $W_{+} \subset W$ of dimension $a$ meets $N$ only in zero, so $a + Q + \dim R \le P + Q + \dim R$, that is $a \le P$; likewise $b \le Q$. Conversely the span of $a$ positive and $b$ negative basis vectors realises $(a,b)$. A spacetime is admitted when $(a,b)$ or $(b,a)$ passes; the two conditions are exchanged by $(p,q) \mapsto (q,p)$, which is the reversal $r \mapsto -r$ along which Premise \claimref{dark-matter.p2-indifference} equates towers, so the admitted set is a union of reversal orbits, as P2 requires (a hypothesis without the sign freedom would admit this world $(1,3)$ and exclude its mirror $(3,1)$, which P2 populates). The check confirms exhaustively that no plane spanned by two vectors with entries in $\{-1, 0, 1\}$ of $\R^{1,4}$ is positive definite, and that coordinate subspaces realise every allowed $(a,b)$.

\emph{(iii).} For $\Lambda > 0$, by (i) and (ii) a causal $(p,q)$ is admitted if and only if $\max(p,q) \le 4$, that is $p + q \le 5$. The causal signatures of dimension two to five are $(1,1)$; $(2,1)$, $(1,2)$; $(3,1)$, $(1,3)$; $(4,1)$, $(1,4)$, with residues $0$; $1$, $7$; $2$, $6$; $3$, $5$: seven signatures, seven distinct residues, residue $4$ absent. A causal signature of residue $r$ has dimension $m + 2$ with $m \equiv \pm r \pmod 8$, so the least is $\min(r, 8-r) + 2$, which is the dimension of the admitted one in each of the seven residues; residue $4$'s least is $6$, $(1,5)$ or $(5,1)$, with $\max(p,q) = 5$. The restriction to causal signatures is the filter (C): without it (W) also admits definite signatures, $(0,4)$ and $(4,0)$ among them, which (C) leaves sterile. For the whole tractor space, $(a,b)$ is admitted if and only if $\min(a,b) \le 2$ and $\max(a,b) \le 4$, which for causal signatures is again $\max(p,q) \le 4$: the same seven, for every $\Lambda$. The check enumerates all causal signatures up to dimension sixteen, compares with an independent enumeration of each residue's lowest causal members, and confirms the sets for $(1,4)$ and $(2,4)$ and their closure under reversal.

\emph{(iv).} For $\Lambda < 0$ the complement is $(2,3)$ and for $\Lambda = 0$ its nondegenerate part is $(1,3)$; in both, by (ii), a causal signature is admitted if and only if $\max(p,q) \le 3$, which removes $(4,1)$ and $(1,4)$, residues $3$ and $5$, and keeps the other five. Checked by the same enumeration.

\emph{(v).} The admitted member of residue $0$ is $(1,1)$ by (iii); there the pseudoscalar $e_{0}e_{1}$ squares to $+1$ and is central in the even algebra (checked), the really chiral case. Theorem \claimref{synthesis.residues-zero-four-real-chirality-anomaly}(iii), conditional on the imported anomaly theorems it names, shows that single-ideal matter at $(1,1)$, massive or massless, has no conserved rotation-covariant stress in the one metric's field and cannot source its field equations. Under (W) residue $0$ has no other admitted member, so with one ideal it holds no tower matter. With both chiral ideals admitted at $(1,1)$ that content is vector-like and anomaly-free (the same theorem's closing remark), so the single-ideal choice there is what remains of the costed clause; residue $4$ needs no clause at all.

\emph{Readings (argued; not part of the check).} (W) is a premise, not a consequence of P1: P1 freezes one $I$, and nothing in the series says that the towers' spacetimes live in its complement. Its support is that the complement is the frozen event's own space, the space whose bivectors are the gauge algebra of the series' MacDowell-Mansouri action (Theorem \claimref{dark-energy.frozen-point-makes-stelle-west-vector}(i): the reduced connection is a connection on $I^{\perp}$), so that (W) reads: every tower's spacetime is a nondegenerate subspace of the space on which gravity's action is a gauge theory; and that Proposition \claimref{dark-matter.nesting} already places this world's algebra inside the residue-$5$ tower's, the tower whose signature the complement has. It was looked for knowing that the window of dimensions up to five gives the unit-ball fit (Proposition \claimref{dark-matter.towers-are-residues-single-clock-filter}(iii) lists that window), so it does not turn $5.364$ into a prediction; it replaces Premise \claimref{dark-matter.p3-nearest-dimension-representatives}, its tie-break included, and the residue-$4$ half of the costed clause by one geometric hypothesis together with the single-clock filter (C); it is recorded as a costed alternative to P3, not adopted. Clause (iv) concerns universes with $\Lambda \le 0$, not a test of this one, since $\Lambda > 0$ was known when (W) was chosen: if (W) were adopted, the five-dimensional towers' share is tied to the sign of $\Lambda$, and a universe with $\Lambda \le 0$ would carry the unit-ball total $1 + 8/3 = 11/3$ under R1 to R3.
